\textbf{Quiz:}
\begin{enumerate}

\item In which scenarios are linear models usually much harder to interpret?
\begin{enumerate}
	\item When the sample size increases.
	\item When their complexity increases (e.g., by adding interaction terms or deeper trees).
	\item When the number of features/main effects increases.
	\item When more assumptions are added.
\end{enumerate}

\item If the model becomes more complex ... 
\begin{enumerate}
	\item ... the interpretation becomes more difficult.
	\item ... the generalization performance improves.
	\item Both
\end{enumerate}

\item In the linear model, the effect and importance of a feature can be inferred from the estimated $\beta$-coefficients. 
\begin{enumerate}
	\item Right
	\item Wrong
\end{enumerate}

\item Which assumptions does the classical linear model impose on the data?
\begin{enumerate}
	\item Normal distribution of $Y \mid X$
	\item Heteroskedasticity
	\item Features dependent on error 
	\item Features uncorrelated
	\item Independent and identically distributed observations
	\item Features are normally distributed
\end{enumerate}

\item What are potential reasons to prefer LASSO over an LM?
\begin{enumerate}
	\item The penalty in LASSO leads to feature selection and sparser models, such that the interpretation is easier. 
	\item The LASSO approach has no additional hyperparameters, both the LM and LASSO need, therefore, no tuning. 
	\item The penalty in LASSO leads to feature selection and sparser models, such that there is less risk of overfitting.
\end{enumerate}

\item Generalized linear models: Why do different link functions exist? 
\begin{enumerate}
	\item To allow for different outcomes, e.g., binary outcomes, ordered categorical outcomes. 
	\item To allow for different interpretations of the coefficients, e.g., logistic vs. probit link, results in different interpretations. 
	\item To allow for different features/covariates, e.g., binary features, ordered categorical features.
\end{enumerate}

\item How are categorical features handled in decision trees?
\begin{enumerate}
	\item They need to be one-hot encoded.
	\item They need to have a natural ordering.
	\item They need no natural ordering.
\end{enumerate}

\item What are the disadvantages of CART? 
\begin{enumerate}
	\item Selection bias towards low-cardinal/discrete features.
	\item There is no proper test whether splitting leads to significant improvements.
	\item They are unstable, such that removing single observations can lead to very different trees.
	\item The shallower they are, the more difficult it is to interpret them.
\end{enumerate}

\end{enumerate}



